% Standalone review draft, 4 November 2026.
% Copy the main subsection and, if useful, the appendix into the manuscript.
% This is a proposed CONTINUUM tip formulation, independent of current code.
% It completes the surface-reduced description with bulk cohesive storage
% in the process zone; it is not a purely notational change to Sections 2.2--2.3.
\documentclass[10pt,a4paper]{article}
\usepackage[margin=24mm]{geometry}
\usepackage[T1]{fontenc}
\usepackage{lmodern,amsmath,amssymb,bm,microtype}
\usepackage{booktabs,tabularx,array}
\usepackage[authoryear,round]{natbib}
\usepackage[colorlinks=true,linkcolor=blue,citecolor=blue,urlcolor=blue]{hyperref}
\setlength{\parindent}{0pt}
\setlength{\parskip}{4pt}
\setlength{\emergencystretch}{2em}
\begin{document}
\begin{center}
{\large\bfseries An energetic treatment of fault tips}\\[3pt]
{\small Proposed main-text subsection and derivation --- 4 November 2026}
\end{center}

\textit{Recommendation.} Keep the physical tip formulation in the main text,
with the supporting intact-limit, tensor, and geometric arguments in an
appendix. The intact limit belongs to the constitutive law, while
irreversibility belongs to the phase equation. Their interaction at a tip
explains their connection but does not justify combining all three subjects
under one heading.

\textit{Modeling choice.} For a propagation model, the mechanical phase driving
must be defined in material that has not yet become a reconstructed fault.
The proposed treatment below achieves this by retaining a bulk cohesive elastic
state in the process zone, and recovering the surface description in the
fault interior. This is one consistent energetic completion, not a universal
consequence of choosing true stress as the primary variable. It requires the
corresponding constitutive extension in Sections 2.2--2.3. The scalar derivation
uses a locally fixed Mode-II frame; the tensor form for changing directions is
given in Appendix~\ref{app:tip_tensor}. Variations are taken at a fixed
instantaneous configuration; global Eulerian transport and geometric work must
be included consistently when that configuration evolves. Constant $G$ and
matched Maxwell relaxation times are assumed below. The final author notes
compare this proposal with the distinct implementations in the repository.

\setcounter{section}{2}
\setcounter{subsection}{3}
% BEGIN MAIN-TEXT SUBSECTION
\subsection{Fault-tip process zone and propagation}
\label{sect:fault_tips}

A structural fault tip separates an existing fault from surrounding intact
material. It is distinct from a rupture front propagating along an existing
fault and from the endpoint of a reconstructed curve. A weakening process zone
connects the slipped fault to intact material; continuing frictional sliding
and the creation of new fracture have distinct energy costs
\citep{FeiChoo2020Shear}. The phase field resolves this transition without
requiring a sharp tip as an internal boundary. Its transverse regularization
length $\ell$ need not equal the along-fault process-zone length.

\paragraph{Constitutive description ahead of the tip.}
The cohesive law is defined in the two-dimensional process zone, including
material ahead of the current reconstructed endpoint. Let
$e_s^{\rm coh}$ be the recoverable cohesive engineering shear strain,
which is dimensionless. For a locally fixed shear direction, select
\begin{equation}
 \psi^{\rm elas}
 =\frac{\bm\tau:\bm\tau}{4G}
    +\frac{G(e_s^{\rm coh})^2}{2h(\phi)},
 \qquad
 \tau_s^{\rm coh}=\frac{G e_s^{\rm coh}}{h(\phi)}.
 \label{eq:tip_bulk_storage}
\end{equation}
The second term is the energy of the additional compliance $h/G$.
Here $e_s^{\rm coh}$ and $\tau_s^{\rm coh}$ are bulk fields. The
surface memory $\bar\tau_s^{\rm coh}$ is their reduced normal-profile
quantity, not an additional copy of the same stored energy.
The mechanical law must likewise supply the admissible shear flow and its
frictional resistance in this process zone; an undefined surface normalization
cannot serve as an intact-material constitutive law.

\paragraph{Phase driving and activation.}
At positive compliance, differentiation at fixed physical elastic states gives
\begin{equation}
 Y:=-\partial_\phi\psi^{\rm elas}
   =\frac{G(e_s^{\rm coh})^2 h'(\phi)}{2h(\phi)^2}
   =\frac{h'(\phi)}{2G}(\tau_s^{\rm coh})^2.
 \label{eq:tip_driving}
\end{equation}
$Y$ is an instantaneous energy-density driving force, not a maximum-history
variable. At $h=0$, the cohesive strain is constrained to zero, while the
cohesive traction remains a reaction determined by the intact mechanical
problem. Its finite, one-sided relaxed limit gives
$Y(0)=h'(0)(\tau_s^{\rm coh})^2/(2G)$, as derived in
Appendix~\ref{app:tip_intact}; no division by zero is required.

With the fracture density introduced in Section~2.3, the phase force is
\begin{equation}
 \mathcal F_\phi
 =\frac{G_c}{c_\alpha\ell}
      \bigl(\alpha'(\phi)-2\ell^2\Delta\phi\bigr)-Y.
 \label{eq:tip_phase_force}
\end{equation}
Rate-independent irreversible evolution satisfies, away from the upper bound,
\begin{equation}
 \dot\phi\geq0,\qquad \mathcal F_\phi\geq0,
 \qquad \dot\phi\,\mathcal F_\phi=0.
 \label{eq:tip_phase_law}
\end{equation}
These equations apply on both sides of the process zone. In a locally uniform
intact region, with $h'(0)>0$, loss of intact stability occurs at
\begin{equation}
 \frac{h'(0)}{2G}(\tau_s^{\rm coh})^2
       =\frac{G_c\alpha'(0)}{c_\alpha\ell}.
 \label{eq:tip_onset}
\end{equation}
This is a local onset condition. Near an existing tip, the full gradient term
in \eqref{eq:tip_phase_force} determines the nonlocal response; the onset
formula is not imposed as an additional pointwise propagation rule.

\paragraph{Advance, direction and reconstruction.}
Damage can advance when the mechanical driving overcomes fracture resistance
in the constrained phase problem. The phase problem therefore
allows variations in the material ahead of the tip, with the full two-dimensional
phase gradient retained. A freely turning fault also requires a constitutive
choice of admissible shear directions in that material. Those directions and
the coupled phase response determine the path; a geometric continuation of the
old tangent is only a prescribed-path restriction. Changes of slip direction
must retain the physical tensor history rather than realign an old scalar
memory (Appendix~\ref{app:tip_tensor}).

The reconstructed curve follows the resulting fault geometry and carries
surface variables where their reduction is valid. A phase threshold used to
locate its numerical endpoint is a representation convention, not a fracture
criterion. No extra phase boundary condition is applied at that endpoint.
The physical lower bound is the transported phase,
$\phi^{n+1}\geq\phi^{n,*}$, and numerical reconstruction does not create or
reset damage. Existing material histories are transferred; newly activated
material follows the stated cohesive and RSF initialization. Fracture energy
is counted once through $\int_\Omega G_c\gamma\,d\bm x$, with no extra
charge for inserting a curve vertex.
% END MAIN-TEXT SUBSECTION

\appendix
\numberwithin{equation}{section}
% BEGIN OPTIONAL APPENDIX
\section{Derivation and interpretation of the tip treatment}

\subsection{Why the driving force must exist before a surface exists}

A bulk phase-field description permits fracture to develop without prescribing
its location as a sharp discontinuity \citep{Bourdin2000Variational}.
A surface reduction alone has a different domain of definition. If mechanical
energy depends on $\phi$ only through normal profiles of an existing curve,
a variation supported wholly outside those profiles changes no mechanical
energy. The fracture-gradient term can smooth a cap there, but does not supply
the missing mechanical driving for activation.

The bulk cohesive state in \eqref{eq:tip_bulk_storage} supplies that driving.
It is physical recoverable deformation; it is not a trial stress or an
operator commutator. A different completion, based on an independently
formulated sharp-front advance law, is possible. The present proposal instead
lets the bulk mechanical and phase laws determine advance, and uses curve
reconstruction to represent the result.

\subsection{The intact limit as a constrained constitutive problem}
\label{app:tip_intact}

The cohesive energy has the extended-valued definition
\begin{equation}
 \psi_s^{\rm elas,coh}(e,h)=
 \begin{cases}
 G e^2/(2h),&h>0,\\
 0,&h=0,\ e=0,\\
 +\infty,&h=0,\ e\ne0.
 \end{cases}
 \label{eq:tip_perspective}
\end{equation}
Here $e=e_s^{\rm coh}$ only within this derivation.
An equivalent mixed form is
\begin{equation}
 \psi_s^{\rm elas,coh}(e,h)
 =\sup_{\tau_s^{\rm coh}}
    \left(\tau_s^{\rm coh}e
          -\frac{h(\tau_s^{\rm coh})^2}{2G}\right).
 \label{eq:tip_mixed}
\end{equation}
For $h>0$, stationarity recovers $e=h\tau_s^{\rm coh}/G$.
For $h=0$, finiteness enforces $e=0$ while the intact mechanical problem
selects its reaction $\tau_s^{\rm coh}$.

To obtain the relaxed one-sided driving, consider an infinitesimal increase
of compliance under that instantaneous conjugate reaction:
\begin{equation}
 \inf_e\left[\frac{Ge^2}{2h}-\tau_s^{\rm coh}e\right]
       =-\frac{h(\tau_s^{\rm coh})^2}{2G}.
 \label{eq:tip_relaxed_energy}
\end{equation}
Its derivative gives $-\partial_\phi\psi_{\rm relaxed}
=h'(0)(\tau_s^{\rm coh})^2/(2G)$ at onset.
This calculation permits the mechanically induced cohesive strain to relax.
It is not the ordinary derivative along the artificially fixed path $e=0$,
which would give zero. The reaction must be obtained from the selected intact
mechanics and friction law; it is not an arbitrary value assigned to a $0/0$
quotient. A smooth regularized RSF law has zero frictional traction at exactly
zero slip rate, whereas a Coulomb sticking law permits a traction interval.
These laws give different reaction selections and must not be interchanged.

Equation~\eqref{eq:tip_onset} has a positive local threshold only when the
chosen $h$ and $\alpha$ supply one. For example, $\alpha'(0)=0$ with
$h'(0)>0$ gives no positive homogeneous elastic onset threshold in this
formulation. This is a model property, not a tip-search parameter.

\subsection{Tensor history and consistency with true-stress mechanics}
\label{app:tip_tensor}

For a fixed shear tensor $\bm{\mathcal S}:\bm{\mathcal S}=1/2$, write
$\bm\varepsilon^{\rm elas,coh}=e_s^{\rm coh}\bm{\mathcal S}$.
For rotating or changing directions retain this tensor itself as the
cohesive elastic history. Select
\begin{equation}
 \psi_s^{\rm elas,coh}
 =\frac{G}{h}\bm\varepsilon^{\rm elas,coh}:
                       \bm\varepsilon^{\rm elas,coh},
 \qquad
 \bm\tau^{\rm coh}=\frac{2G}{h}\bm\varepsilon^{\rm elas,coh}.
 \label{eq:tip_tensor_storage}
\end{equation}
The scalar traction is $\bm\tau^{\rm coh}:\bm{\mathcal S}$.
For a single fixed shear direction, \eqref{eq:tip_tensor_storage} reduces
exactly to \eqref{eq:tip_bulk_storage}.
For general tensor memory, its phase driving is instead
\begin{equation}
 Y=\frac{G h'}{h^2}\bm\varepsilon^{\rm elas,coh}:
                         \bm\varepsilon^{\rm elas,coh}
   =\frac{h'}{4G}\bm\tau^{\rm coh}:\bm\tau^{\rm coh}.
 \label{eq:tip_tensor_driving}
\end{equation}
It cannot generally be replaced by the square of one resolved shear traction.
This tensor storage is an explicit constitutive extension of the single-direction
model; its admissible strain space and frictional direction law must be specified.

With the bulk fault-deformation rate $\dot{\bm\varepsilon}^{\rm fault}$,
one consistent matched-Maxwell choice is
\begin{equation}
 \frac{\mathring{\bm\tau}}{2G}+\frac{\bm\tau}{2\eta}
   =\dot{\bm\varepsilon}-\dot{\bm\varepsilon}^{\rm fault},
 \qquad
 (\bm\varepsilon^{\rm elas,coh})^\circ
     +\frac{G}{\eta}\bm\varepsilon^{\rm elas,coh}
   =\dot{\bm\varepsilon}^{\rm fault}.
 \label{eq:tip_tensor_rates}
\end{equation}
The fault-deformation rate is defined throughout the process zone. Its
interior approximation may be $\chi\widetilde V\widetilde{\bm{\mathcal S}}$;
that approximation is not used to define an intact $0/0$ normalization.
The physical slip/flow law must make
$(\bm\tau-\bm\tau^{\rm coh}):\dot{\bm\varepsilon}^{\rm fault}$
its nonnegative frictional and radiation loss, with a dimensionally consistent
relation between a bulk deformation rate and any interface slip rate.

Contracting \eqref{eq:tip_tensor_rates} with the conjugate stresses gives the
local stored-energy rate
\begin{equation}
 \begin{split}
 \dot\psi^{\rm elas}={}&\bm\tau:\dot{\bm\varepsilon}
 -\frac{\bm\tau:\bm\tau}{2\eta}
 -\frac{h\bm\tau^{\rm coh}:\bm\tau^{\rm coh}}{2\eta}\\
 &-(\bm\tau-\bm\tau^{\rm coh}):\dot{\bm\varepsilon}^{\rm fault}
 -Y\dot\phi.
 \end{split}
 \label{eq:tip_local_power}
\end{equation}
The stress-spin contractions vanish. Thus $Y\dot\phi$ is the mechanical
power made available to phase evolution, while bulk viscosity, cohesive
viscosity, friction and radiation remain separately accounted losses.
No total displacement history, auxiliary trial stress or localization
commutator enters this derivation.

\subsection{Connection to the surface cohesive energy}

For a locally straight interior profile with uniform
$\tau_s^{\rm coh}$ across it,
\begin{equation}
 \bar\tau_s^{\rm coh}:=G\int e_s^{\rm coh}\,d\zeta
       =\bar h_\xi\tau_s^{\rm coh},
 \qquad
 \int\frac{G(e_s^{\rm coh})^2}{2h}\,d\zeta
       =\frac{(\bar\tau_s^{\rm coh})^2}{2G\bar h_\xi}.
 \label{eq:tip_surface_reduction}
\end{equation}
This recovers the interior surface energy. It uses a locally flat profile,
line measure $d\zeta$, and uniform cohesive traction; general curved/nonuniform
profiles require their actual measures and retain additional transverse state.
Indeed, for $\bar h_\xi=\int h\,d\zeta>0$, the exact flat-profile identity is
\begin{equation}
 \int\frac{G(e_s^{\rm coh})^2}{2h}\,d\zeta
 -\frac{(\bar\tau_s^{\rm coh})^2}{2G\bar h_\xi}
 =\int\frac{h}{2G}
 \left(\tau_s^{\rm coh}-\frac{\bar\tau_s^{\rm coh}}{\bar h_\xi}\right)^2d\zeta.
 \label{eq:tip_transverse_variance}
\end{equation}
Thus a single surface moment loses recoverable energy when the transverse
traction is nonuniform; recovering it is a constitutive approximation, not
just a change of notation.
The bulk and surface expressions describe the same cohesive storage and are
not added over the same region. Near a tip, where longitudinal variation and
intact activation matter, the bulk description remains authoritative.

\subsection{Variational growth and its sharp-front interpretation}

At the time-step endpoint, let $\phi^{n,*}$ be the transported accepted lower bound.
The weak phase law is
\begin{equation}
 \begin{gathered}
 \phi^{n,*}\leq\phi^{n+1}\leq1,\\
 \int_\Omega\left[
 \left(\frac{G_c\alpha'(\phi^{n+1})}{c_\alpha\ell}-Y^{n+1}\right)
                  (\widehat\phi-\phi^{n+1})
 +\frac{2G_c\ell}{c_\alpha}\nabla\phi^{n+1}\cdot
                    \nabla(\widehat\phi-\phi^{n+1})\right]d\bm x\geq0
 \end{gathered}
 \label{eq:tip_vi}
\end{equation}
for every admissible $\widehat\phi$. Mechanical and cohesive states are
updated consistently with this phase problem. A negative first variation in
an admissible direction ahead of the current tip requires a phase update;
a chosen solver must also address nonconvex branch selection. Stationarity
alone does not prove global stability.

For a single sharp front in a quasistatic rate-independent limit, define
$\mathcal G_{\rm II}$ as the net mechanical energy available per unit new
fault area after the appropriate sliding/viscous losses and transport/geometry
work are accounted for. If $c_{\rm tip}$ is the nonnegative advance speed
relative to the material, fracture power is $G_c c_{\rm tip}$ per unit
out-of-plane thickness. Energetic stability and balance give
\begin{equation}
 c_{\rm tip}\geq0,\qquad
 G_c-\mathcal G_{\rm II}\geq0,\qquad
 c_{\rm tip}(G_c-\mathcal G_{\rm II})=0.
 \label{eq:tip_griffith}
\end{equation}
This is the sharp-front interpretation of an appropriately calibrated and
convergent phase formulation, not an extra condition to impose at every
phase-field tip. At finite $\ell$, the diffuse fracture functional, rather
than $G_c$ times a threshold-defined curve length, is the actual energy.
An elastic traction-free $J$-integral is not automatically
$\mathcal G_{\rm II}$ for a frictional viscoelastic fault.

\subsection{Geometry and endpoint conditions}

For a straight fault $\bm f(\xi)=(\xi,0)$ ending at $\xi=0$, normal
coordinates $\bm f(\xi)+\zeta\hat{\bm n}=(\xi,\zeta)$ with $\xi\leq0$
cover the strip behind the tip. A two-dimensional cap ahead of it projects
onto the single endpoint, not onto a family of normal fibers. Assigning that
cap the endpoint slip value would create a concentrated endpoint contribution
to the work adjoint. The bulk tip formulation avoids treating such a cap as
an ordinary interior surface profile. Boundary completion at a domain edge
remains distinct from physical tip advance.

The tip velocity and slip rate are also distinct. Writing its position as
$\bm f(\xi_{\rm tip}(t),t)$ gives
\begin{equation}
 \frac{d}{dt}\bm f(\xi_{\rm tip}(t),t)
 =\partial_t\bm f+\partial_\xi\bm f\,\dot\xi_{\rm tip}.
 \label{eq:tip_motion}
\end{equation}
For illustration, on a stationary straight material path introduce sharp slip
$U(\xi,t)$ with $V=\partial_tU|_\xi$. If the two sides meet at an advancing
front, $U(\xi_{\rm tip}(t),t)=0$ implies, wherever the derivatives exist,
\begin{equation}
 V+\dot\xi_{\rm tip}\partial_\xi U=0.
 \label{eq:tip_compatibility}
\end{equation}
It does not generally imply $V=0$ at a moving front. Singular sharp-tip
fields require a limiting interpretation. Zero rate may be appropriate for
a stationary tied endpoint or another specifically derived process-zone
condition; it is not a universal condition at the last reconstructed vertex.
$U$ here explains compatibility and need not be stored in the rate formulation.

The phase equation has no internal tip boundary. Its outer boundary may carry
natural zero microtraction or prescribed phase data. Tip position, candidate
spacing and a geometric threshold are numerical diagnostics/settings; physical
activation follows the coupled law. Verify growth threshold, path, advance and
work balance under bulk/surface/time refinement at fixed physical $\ell$.
Varying $\ell$ is a separate model/regularization study. Transferring an existing
material point preserves its history; genuinely new material requires a stated
initial cohesive state and positive RSF age, not a reset based on a new index.
% END OPTIONAL APPENDIX

\begin{thebibliography}{2}
\bibitem[Bourdin et~al.(2000)]{Bourdin2000Variational}
B. Bourdin, G. A. Francfort, and J.-J. Marigo (2000).
Numerical experiments in revisited brittle fracture.
\textit{Journal of the Mechanics and Physics of Solids}, 48, 797--826.
\href{https://doi.org/10.1016/S0022-5096(99)00028-9}
{doi:10.1016/S0022-5096(99)00028-9}.
\bibitem[Fei and Choo(2020)]{FeiChoo2020Shear}
F. Fei and J. Choo (2020).
A phase-field model of frictional shear fracture in geologic materials.
\textit{Computer Methods in Applied Mechanics and Engineering}, 369, 113265.
\href{https://doi.org/10.1016/j.cma.2020.113265}
{doi:10.1016/j.cma.2020.113265}.
\end{thebibliography}

\clearpage
% BEGIN AUTHOR NOTES -- do not paste this implementation audit into the theory section.
\begingroup
\small
\setlength{\parskip}{3pt}
\section*{Author notes: continuum target versus current implementation}

This comparison concerns branch \texttt{pf-rsf-particle-free} at commit
\href{https://github.com/YiminJin/aspect/tree/69e8a10a93a08718d31af961197899e320657c04}
{\texttt{69e8a10a93a0}}, reviewed on 4 November 2026. It is based on source
inspection and the repository's recorded tests; those simulations were not rerun
for this review. The repository contains several distinct formulations. An
absent capability in the older surface-reduced route must not be attributed to
all of them.

\renewcommand{\arraystretch}{1.08}
\begin{tabularx}{\linewidth}{@{}>{\raggedright\arraybackslash}p{.18\linewidth}
 >{\raggedright\arraybackslash}p{.31\linewidth}
 >{\raggedright\arraybackslash}X@{}}
\toprule
Topic & Continuum target in this draft & Current implementation and difference\\
\midrule
Driving ahead of a tip & A bulk cohesive reaction and phase force exist before
there is a surface record. & Already implemented in the separate bulk cohesive
route, including its $h=0$ reaction. The older surface-profile route has zero
mechanical phase derivative strictly ahead of its normal strip. [1,2]\\
\addlinespace
Turning and memory & Retain cohesive tensor history when the active direction
changes; solve the full spatial phase problem. & Native fixed/radial tensor
bulk modes already retain this history and its consistent phase coupling.
Reported tests show connected activation in initially intact material, but
not converged general turning, branching or developed crack propagation. [3]\\
\addlinespace
Friction and RSF & Specify a process-zone flow law and its relation to physical
slip rate, with consistent friction/radiation work. & The evolving bulk mode
uses prescribed constant resistance; the pressure-dependent Coulomb option
requires fixed phase. Bulk RSF is explicitly unsupported. A separate interface
pilot has RSF and segment admission, with a different physical damage variable
and fracture functional. [2,4]\\
\addlinespace
Motion and growth & Material transport, objective history, phase irreversibility
and new activation must coexist consistently. & Restricted moving mature-fault
experiments transport history/profile with growth disabled. The bulk cohesive
growth route excludes transport and moving mesh; general deformed cohesive
evolution remains unsupported. [2,5]\\
\addlinespace
Energy and accuracy & Count physical fracture once and establish growth/path/work
convergence independently of algebraic solve accuracy. & Bulk and interface
routes already have separate storage/dissipation accounting and accepted-state
publication. Their exact discrete balances do not establish propagation
accuracy. The interface's later localized-tip and fracture-relative physical
energy screens retain failures. [3,4,6]\\
\bottomrule
\end{tabularx}

\paragraph{What this means for the proposed section.}
The bulk cohesive route is the closest existing implementation of the
continuum construction above. Its core phase driving and intact limit do
\emph{not} need to be invented again. The remaining difference is a complete,
qualified coupling to RSF, general material motion, and the intended fault
geometry representation. Conversely, the newer cohesive-interface pilot is
an alternative physical formulation, not merely a more efficient discretization
of this bulk phase equation: it charges $G_c\int_\Gamma a\,ds$, while this
draft charges $\int_\Omega G_c\gamma(\phi,\nabla\phi)\,d\bm x$.
Both may be useful, but their damage variables and energy inventories cannot
be identified without a separate derivation.

The code also retains auxiliary stress representations. Changing their storage
to direct $\bm\tau$ is a separate representation/lifecycle task; it does not
by itself supply tip physics, RSF, or transport. No commutator is needed in
the direct physical constitutive equations given here.

\paragraph{Source anchors.}
\begin{enumerate}
\setlength{\itemsep}{1pt}
\setlength{\parskip}{0pt}
\item \href{https://github.com/YiminJin/aspect/blob/69e8a10a93a08718d31af961197899e320657c04/doc/reconstructed_fault/particle_free/P4A_TIP_REPORT.md}{P4a finite-tip report}:
strictly-ahead derivatives and the restricted older surface route.
\item \href{https://github.com/YiminJin/aspect/blob/69e8a10a93a08718d31af961197899e320657c04/source/material_model/phase_field_fault/bulk_cohesive.cc}{Bulk scalar return},
\href{https://github.com/YiminJin/aspect/blob/69e8a10a93a08718d31af961197899e320657c04/source/material_model/phase_field_fault/bulk_cohesive_tensor.cc}{tensor return},
\href{https://github.com/YiminJin/aspect/blob/69e8a10a93a08718d31af961197899e320657c04/source/simulator/solver/bulk_cohesive_phase.cc}{coupled phase solver}, and
\href{https://github.com/YiminJin/aspect/blob/69e8a10a93a08718d31af961197899e320657c04/source/material_model/phase_field_fault.cc}{material-model admission guards}.
\item \href{https://github.com/YiminJin/aspect/blob/69e8a10a93a08718d31af961197899e320657c04/doc/reconstructed_fault/particle_free/P5B_NATIVE_REPORT.md}{P5b native tensor integration} and
\href{https://github.com/YiminJin/aspect/blob/69e8a10a93a08718d31af961197899e320657c04/doc/reconstructed_fault/particle_free/P5C_CONTINUATION_REPORT.md}{P5c connected process-zone extension}:
implemented histories, measured activation, and remaining convergence limits.
\item \href{https://github.com/YiminJin/aspect/blob/69e8a10a93a08718d31af961197899e320657c04/doc/reconstructed_fault/particle_free/COHESIVE_RSF_NATIVE_REPORT.md}{Native cohesive--RSF report} and
\href{https://github.com/YiminJin/aspect/blob/69e8a10a93a08718d31af961197899e320657c04/doc/reconstructed_fault/particle_free/IMPLEMENTATION_MAPS.md}{implementation/energy inventory}:
surface damage, formation, histories and alternative fracture accounting.
\item \href{https://github.com/YiminJin/aspect/blob/69e8a10a93a08718d31af961197899e320657c04/doc/reconstructed_fault/current_design.md}{Current design, Stage C material-step sections}:
restricted moving mature model and deferred cohesive geometry work.
\item \href{https://github.com/YiminJin/aspect/blob/69e8a10a93a08718d31af961197899e320657c04/doc/reconstructed_fault/particle_free/STAGE_C_COHESIVE_DIAGNOSIS.md}{Stage C cohesive diagnosis}:
physical energy remainder and localized-tip nonlinear failure.
\end{enumerate}
\endgroup
% END AUTHOR NOTES
\end{document}
